\section{Setup and Notation}

We fix a time interval \([0,T)\) and a classical solution \(u\) on
\([0,T)\times\mathbb{R}^3\) arising from smooth compactly supported
divergence-free data. The background PDE framework is the standard one
developed from Leray and Hopf weak solutions and organized in the monographs of
Ladyzhenskaya and Temam \cite{Leray1934,Hopf1951,Ladyzhenskaya1969,Temam2001}.
We write \(P_{\ge N}\) for a dyadic high-frequency projector and
\(\Pi_{\ge N}(t)\) for the corresponding high-frequency flux quantity used in
the scale-barrier hypothesis.

\begin{definition}[Scale barrier]
A scale barrier is a classical estimate on \(\Pi_{\ge N}(t)\) that blocks unbounded transfer of activity into dyadic scales associated with singularity formation.
\end{definition}

\begin{definition}[Monotone functional]
The monotone functional is the classical control quantity
\[
Q(t)=\|u(t)\|_{L^2_x}^2+\|\nabla u(t)\|_{L^2_x}^2,
\]
used by the original source stack to track the global energy/enstrophy layer on the same equation surface as the later compactness and continuation arguments.
\end{definition}

\begin{definition}[Nonlinear compactness upgrade]
A nonlinear compactness upgrade is a classical approximation scheme \(u^{(n)}\to u\) whose convergence is strong enough to pass the nonlinear interaction \(u^{(n)}\otimes u^{(n)}\to u\otimes u\) at the theorem level.
\end{definition}

\begin{definition}[Classical gradient control]
Classical gradient control is derivative-scale control, here measured by \(\|\nabla u(t)\|_{L^2(\mathbb{R}^3)}\), derived from the same classical dynamics and not from added regularization or geometry.
\end{definition}

The proof architecture below is built around these four bridges plus the non-dependence condition on modified equations.
